\section{Methods}
\label{sec:methods}

We formulate snatch phase identification as supervised temporal labeling from fixed pose features.
This section documents only the \emph{verified} baseline pipeline; subsections for additional benchmark architectures are reserved with explicit \textsc{TODO} markers.

\subsection{Problem formulation}
\label{sec:methods:problem}

Let each video $v$ yield a pose sequence $\mathbf{X}^{(v)} \in \mathbb{R}^{T \times \featuredim}$ after feature extraction, where $T$ is the number of frames.
Frame-wise ground truth is $y^{(v)}_t \in \{0,\ldots,7\}$ with phase~0 denoting \texttt{unlabeled}.

The baseline consumes sliding windows $\mathbf{w}_i \in \mathbb{R}^{\windowsize \times \featuredim}$ and predicts a single class label $\hat{y}_i$ for window~$i$, supervised by the center-frame label $y_{t(i)}$.
Benchmark extensions will additionally produce frame-wise labels $\hat{y}^{(v)}_t$ and contiguous predicted segments for segment-level metrics (\Cref{sec:protocol:metrics}).

\subsection{End-to-end pipeline}
\label{sec:methods:pipeline}

\Cref{fig:pipeline} summarizes the verified processing chain:
(i)~pose keypoints per frame,
(ii)~alignment with frame labels,
(iii)~sliding-window tensor construction,
(iv)~train-only feature standardization,
(v)~\acs{LSTM} training with athlete-level splits,
(vi)~window-level evaluation on held-out athletes.

\begin{figure}[t]
  \centering
  \todofig{End-to-end pipeline: raw video (optional) $\rightarrow$ MediaPipe keypoints $\rightarrow$ window tensor dataset $\rightarrow$ LSTM classifier $\rightarrow$ window-level metrics. Regenerate from \texttt{docs/figures\_plan.md} F2.}
  \caption{Verified SnatchPhaseBench baseline pipeline. Pose CSV files in the current artifact replace on-the-fly extraction when raw video or the MediaPipe task binary is unavailable.}
  \label{fig:pipeline}
\end{figure}

\subsection{Preprocessing}
\label{sec:methods:preprocessing}

Preprocessing follows the frozen implementation ported to the canonical repository (\texttt{dataset\_builder.py}; see \texttt{docs/FROZEN\_BASELINE.md}).

\paragraph{Alignment.}
Frame labels from \texttt{master\_frame\_labels.csv} are merged with keypoint rows on $(\text{video\_relpath}, \text{frame})$.

\paragraph{Missing pose values.}
Per-coordinate linear interpolation across time is applied, followed by forward and backward filling for remaining gaps.

\paragraph{Window construction.}
For each video, all windows of length $\windowsize = 31$ with stride~1 are extracted.
The supervised label is the phase ID at the center frame (index 15 within the window).
Windows whose center label is \texttt{unlabeled} (ID~0) are removed.

\paragraph{Feature tensor.}
Each window is stored as a float32 array of shape $(\windowsize, \featuredim)$ with $\featuredim = 99$ from $(x,y,z)$ coordinates of \numlandmarks{} landmarks.
Visibility channels are excluded in the baseline configuration.

\paragraph{Standardization.}
Before training, features are standardized using mean and variance computed on \emph{training-split windows only}, then applied to validation and test windows.

\subsection{Baseline model: LSTM classifier}
\label{sec:methods:lstm}

The verified baseline is a single-layer \acs{LSTM} followed by dropout and a linear classification head (\texttt{lstm\_trainer.py} / \texttt{LSTMClassifier}).
Given window input $\mathbf{w} \in \mathbb{R}^{\windowsize \times \featuredim}$, the model outputs logits over \numclasses{} training phases and is trained with cross-entropy loss using class-frequency weights computed on the training split.

\Cref{tab:lstm_hyperparams} lists hyperparameters frozen in \texttt{configs/baseline\_lstm.yaml}.

\begin{table}[t]
  \centering
  \caption{Verified \acs{LSTM} baseline hyperparameters (frozen configuration).}
  \label{tab:lstm_hyperparams}
  \begin{tabular}{ll}
    \toprule
    Setting & Value \\
    \midrule
    Hidden size & 128 \\
    \acs{LSTM} layers & 1 \\
    Dropout & 0.2 \\
    Optimizer & Adam ($\mathrm{lr}=10^{-3}$, weight decay $10^{-4}$) \\
    Batch size & 64 \\
    Max epochs & 40 \\
    Early stopping & Validation macro-F1, patience 8 \\
    Class weighting & Enabled (inverse frequency) \\
    Standardization & Train split only \\
    Random seed & 42 \\
    \bottomrule
  \end{tabular}
\end{table}

The final checkpoint selection criterion monitors validation macro-F1.
This window-level classifier constitutes the initial baseline for SnatchPhaseBench; it is not claimed to be state of the art.

\subsection{Benchmark models (planned)}
\label{sec:methods:benchmark}

The software registry supports plugging additional temporal models without modifying the frozen baseline modules.
The following subsections are placeholders for Phase~3 experiments.
\emph{No architectures below have been trained or evaluated in the canonical benchmark yet.}

\subsubsection{GRU baseline}
\label{sec:methods:gru}
\textsc{TODO:} Describe gated recurrent baseline with matched parameter budget once implemented.

\subsubsection{Temporal convolutional network (TCN)}
\label{sec:methods:tcn}
\textsc{TODO:} Describe dilated causal TCN head producing frame-wise logits; cite verified \acs{TCN} literature when added.

\subsubsection{Multi-stage TCN (MS-TCN)}
\label{sec:methods:mstcn}
\textsc{TODO:} Describe multi-stage refinement architecture for dense segmentation.

\subsubsection{Spatial--temporal graph convolutional network (ST-GCN)}
\label{sec:methods:stgcn}
\textsc{TODO:} Define skeleton graph over MediaPipe landmarks; specify input window and readout for phase labels.

\subsubsection{Transformer encoder}
\label{sec:methods:transformer}
\textsc{TODO:} Specify tokenization of pose frames, positional encoding, and prediction head.

\subsection{Window construction illustration}
\label{sec:methods:windows}

\begin{figure}[t]
  \centering
  \todofig{Sliding window of size 31 with stride 1; center-frame label assignment; unlabeled centers discarded. See \texttt{docs/figures\_plan.md} F3.}
  \caption{Sliding-window labeling protocol (verified). Adjacent windows overlap substantially when stride~$=1$, inducing temporal dependence among training and evaluation samples (\Cref{sec:protocol:overlap}).}
  \label{fig:window_construction}
\end{figure}
